\subsection{Results on the extended ranges}

\begin{table}[ht]\centering\small
% generated by scripts/c128/summarise_ext.py -- do not edit
% src: data/msc_rigorous_certified/SUMMARY_EXT.json, key modes (from data/msc_rigorous_certified/c128_*.jsonl, data/msc_rigorous_certified/xcheck_limb_*.jsonl, data/msc_rigorous_certified/xcheck_packed_*.jsonl; times: sums of the field seconds, wall-clock)
\begin{tabular}{cc|p{3.2cm}|r|r|p{3.1cm}|p{2.5cm}|r}
$d$ & $q$ & certified $N$ (C engine) & number & new & re-derived by \texttt{fplimb} & by \texttt{verify\_packed} & time (s) \\ \hline
1 & $4/5$ & $2$--$1500$, $2000$ & 1500 & 899 & 650, 700, 750, 850, 900, 950, 1200, 1500 & 700, 900 & 5310 \\
1 & $1/2$ & $2$--$1000$ & 999 & 700 & 500, 900 & 400, 700 & 1343 \\
1 & $1$ & $2$--$400$ & 399 & 200 & 300 & 300 & 175 \\
2 & $4/5$ & $2$--$301$, $350, 401, 450, 500, 600$ & 305 & 144 & 210, 220, 230, 250 & 170, 190 & 9351 \\
2 & $1/2$ & $2$--$200$ & 199 & 120 & 120, 190 & 100, 170 & 1228 \\
2 & $1$ & $2$--$160$, $180, 200$ & 161 & 82 & 120, 136, 149 & 86, 92, 136, 149 & 275 \\
3 & $4/5$ & $2$--$120$, $130, 140, 150, 160$ & 123 & 64 & 62, 64, 66, 70, 75 & -- & 9986 \\
3 & $1/2$ & $2$--$80$ & 79 & 48 & 40, 64 & 36 & 865 \\
3 & $1$ & $2$--$60$ & 59 & 28 & 40 & 40 & 90 \\
\end{tabular}

\caption{Ranges of $N$ covered by the certificates of the C engine (\texttt{certs/c128\_*.jsonl}). ``New'': the number of values of $N$ not covered by Table~\ref{tab:ranges}. Columns 6 and 7: the new values of $N$ whose certificates were re-derived by the programs of Section~\ref{sec:lemmas} (identical $t_0$, $t^*$, $T$ and verdicts, intersecting enclosures). Time: total wall-clock time in seconds of the C engine and of the Python post-processing, on a heavily loaded machine.}\label{tab:rangesX}
\end{table}

For $d\in\{1,2,3\}$ and $q\in\{4/5,1/2,1\}$ let $\mathcal N_{\rm C}(d,q)\supset\mathcal N(d,q)$ be the set of $N$ in the third column of Table~\ref{tab:rangesX}:
\begin{itemize}
\item[$q=4/5$:] $d=1$: $2\le N\le\XNmaxOneF$\XAlsoOneF;\quad $d=2$: $2\le N\le\XNmaxTwoF$\XAlsoTwoF;\quad $d=3$: $2\le N\le\XNmaxThreeF$\XAlsoThreeF.
\item[$q=1/2$:] $d=1$: $2\le N\le\XNmaxOneH$\XAlsoOneH;\quad $d=2$: $2\le N\le\XNmaxTwoH$\XAlsoTwoH;\quad $d=3$: $2\le N\le\XNmaxThreeH$\XAlsoThreeH.
\item[$q=1$:] $d=1$: $2\le N\le\XNmaxOneU$\XAlsoOneU;\quad $d=2$: $2\le N\le\XNmaxTwoU$\XAlsoTwoU;\quad $d=3$: $2\le N\le\XNmaxThreeU$\XAlsoThreeU.
\end{itemize}

\begin{theorem}[modes and unimodality on the extended ranges; computer-assisted, trusted base (T5)]\label{thm:modesX}
Let $d\in\{1,2,3\}$, $q\in\{4/5,\,1/2\}$ and $N\in\mathcal N_{\rm C}(d,q)$, and suppose that $(d,q,N)$ is not one of the three exceptional triples of Theorem~\ref{thm:modes}.
\begin{itemize}
\item[(a)] $f$ is strictly unimodal in the sense of Definition~\ref{def:unimodal}: $f(t)=0$ for $t<d(N-1)$, $f(t)<f(t+1)$ for $d(N-1)-1\le t<t^*$ and $f(t)>f(t+1)$ for all $t\ge t^*$, where $t^*=t^*(d,N,q)$ is the integer in the field \texttt{mode} of line $N$ of \texttt{certs/c128\_d$d$\_q4-5.jsonl} (resp.\ \texttt{\dots\_q1-2.jsonl}). Table~\ref{tab:mX} lists a selection, Appendix~\ref{sec:allmodes} all new values for $d=2,3$, and Theorem~\ref{thm:1dlawX} a formula for $d=1$.
\item[(b)] $v_{t^*+1}<v_{t^*}$ on $\Omega'$ (the tail time is $T=t^*$), and $f(t^*+s)\le\gamma^sf(t^*)$ for all $s\ge0$ with a number $\gamma=\gamma(d,N,q)<1$ such that $1-\gamma\ge$ the field \texttt{one\_minus\_gamma\_lower}; over the ranges, $1-\gamma\ge\XGamOneF$, $\XGamTwoF$, $\XGamThreeF$ for $d=1,2,3$ at $q=4/5$, and $1-\gamma\ge\XGamOneH$, $\XGamTwoH$, $\XGamThreeH$ at $q=1/2$.
\end{itemize}
For the three exceptional triples the C certificate reports exactly one undecided relation, in agreement with the exact ties of Theorem~\ref{thm:modes}(d). No further exceptional case occurs in $\mathcal N_{\rm C}(d,q)$.
% src: data/msc_rigorous_certified/SUMMARY_EXT.json (key modes: min_one_minus_gamma, not_certified_N, non_unimodal_N); fields mode, T_tail, one_minus_gamma_lower, undecided_adjacent_pairs of data/msc_rigorous_certified/c128_d*_q4-5.jsonl and data/msc_rigorous_certified/c128_d*_q1-2.jsonl
\end{theorem}
\begin{proof}
For each $(d,q,N)$ the certificate was produced by \texttt{fp128} and \texttt{c128\_post.py}. By Proposition~\ref{prop:soundC} its data satisfy the hypotheses of Proposition~\ref{prop:sound}, whose parts (i)--(iii) give (a) and (b) from the recorded verdicts (mode certified, strictly unimodal, no undecided relation, $T=t^*$); the first passage time $d(N-1)$ is Lemma~\ref{lem:support}. The script \texttt{scripts/c128/audit\_ext.py} re-checks every assertion against the certificate files (\texttt{certs/AUDIT\_EXT.json}).
\end{proof}

\begin{theorem}[modes and multimodality for $q=1$ on the extended ranges; computer-assisted, trusted base (T5)]\label{thm:modesUX}
Let $q=1$, $d\in\{1,2,3\}$ and $N\in\mathcal N_{\rm C}(d,1)$. The smallest maximiser of $f$ is the integer in the field \texttt{mode} of \texttt{certs/c128\_d$d$\_q1-1.jsonl}, and:
% generated by scripts/c128/summarise_ext.py -- do not edit
% src: data/msc_rigorous_certified/c128_d*_q1-1.jsonl, data/msc_rigorous_certified/packedexact_d*_q1-1.jsonl, data/msc_rigorous_certified/xcheck_packed_d2_q1-1.jsonl; exact ties re-derived by code/msc_rigorous_certified/verify_exact.py
\begin{itemize}
\item[$d=1$:] the maximiser is unique except for $N\in\{5, 6, 8, 9\}$, where the maximum is attained at exactly two times ($N=5$: $t\in\{4, 6\}$; $N=6$: $t\in\{7, 9\}$; $N=8$: $t\in\{15, 17\}$; $N=9$: $t\in\{20, 22\}$; exact arithmetic, Theorem~\ref{thm:modesU}); $f$ is strictly unimodal for $N\in{}$\{2\}; $f$ has at least two strict local maxima, hence is \emph{not} unimodal, for $N\in{}$\{4--400\} (the largest certified number of strict local maxima is 164054); for $N=3$, $f$ is unimodal in the weak sense but not strictly: $f(3)=f(4)$; equal consecutive non-zero values occur only for $N=3$ ($f(3)=f(4)$) and $N=4$ ($f(8)=f(9)$) (exact arithmetic, Theorem~\ref{thm:modesU}): for every other pair $(N,t)$ with $N$ in the range and $t\ge d(N-1)-1$, $f(t)\ne f(t+1)$; the tail time $T$ exceeds $t^*$ for 398 of the 399 values of $N$ (largest ratio $T/t^*=6.19$).
\item[$d=2$:] the maximiser is unique except for $N\in\{2\}$, where the maximum is attained at exactly two times ($N=2$: $t\in\{2, 3\}$; exact arithmetic, Theorem~\ref{thm:modesU}); $f$ is strictly unimodal for $N\in{}$\{3--8, 10, 12, 14, 16, 19--20, 22--23, 25, 27--37, 39--50, 52--60, 62--85, 87--91, 93--135, 137--148, 150--160, 180, 200\}; $f$ has at least two strict local maxima, hence is \emph{not} unimodal, for $N\in{}$\{9, 11, 13, 15, 17--18, 21, 24, 26, 38, 51, 61, 86, 92, 136, 149\} (the largest certified number of strict local maxima is 2); for $N=2$, $f$ is unimodal in the weak sense but not strictly: $f(2)=f(3)$; equal consecutive non-zero values occur only for $N=2$ ($f(2)=f(3)$) (exact arithmetic, Theorem~\ref{thm:modesU}): for every other pair $(N,t)$ with $N$ in the range and $t\ge d(N-1)-1$, $f(t)\ne f(t+1)$; the tail time $T$ exceeds $t^*$ for 71 of the 161 values of $N$ (largest ratio $T/t^*=1.50$).
\item[$d=3$:] the maximiser is unique for every $N$ of the range; $f$ is strictly unimodal for $N\in{}$\{2--60\}; $f(t)\ne f(t+1)$ for every $N$ of the range and every $t\ge d(N-1)-1$.
\end{itemize}

\end{theorem}
\begin{proof}
Propositions~\ref{prop:soundC} and \ref{prop:sound}(i)--(iv). Ties are not decided by the C engine; the cases with a tie are those of Theorem~\ref{thm:modesU}, settled there in exact arithmetic. The four two-dimensional boxes with two strict local maxima that are not in Theorem~\ref{thm:modesU}, $N=86,92,136,149$, were re-derived with Python integers only (\texttt{verify\_packed.py}, Table~\ref{tab:rangesX}), so these four statements do not depend on (T5).
% src: data/msc_rigorous_certified/xcheck_packed_d2_q1-1.jsonl (N = 86, 92, 136, 149); data/msc_rigorous_certified/packedexact_d*_q1-1.jsonl
\end{proof}

\begin{table}[ht]\centering\footnotesize
% generated by scripts/c128/summarise_ext.py -- do not edit
% src: data/msc_rigorous_certified/closedform_c128_d*_q4-5.jsonl (t*, tau_N, t_cf), data/msc_rigorous_certified/c128_d*_q1-2.jsonl, data/msc_rigorous_certified/c128_d*_q1-1.jsonl
\begin{tabular}{c|r|r|r|r|r|r|r|r}
$d$ & $N$ & $t^*$ ($q=4/5$) & $\tau_N=q\,\mathbb E T$ & $t^*/\mathbb E T$ & $t_{\rm cf}$ & $(t_{\rm cf}-t^*)/t^*$ & $t^*$ ($q=1/2$) & $t^*$ ($q=1$) \\ \hline
1 & 600 & 149727 & 359400.0000 & 0.333282 & 173420.722 & $+0.15825$ & 239564 & -- \\
1 & 700 & 203844 & 489300.0000 & 0.333283 & 236078.917 & $+0.15814$ & 326151 & -- \\
1 & 800 & 266294 & 639200.0000 & 0.333284 & 308381.324 & $+0.15805$ & 426070 & -- \\
1 & 900 & 337075 & 809100.0000 & 0.333284 & 390327.944 & $+0.15799$ & 539320 & -- \\
1 & 1000 & 416188 & 999000.0000 & 0.333284 & 481918.777 & $+0.15794$ & 665901 & -- \\
1 & 1200 & 599411 & 1438800.0000 & 0.333284 & 694033.080 & $+0.15786$ & -- & -- \\
1 & 1500 & 936737 & 2248500.0000 & 0.333284 & 1084536.129 & $+0.15778$ & -- & -- \\
1 & 2000 & 1665588 & 3998000.0000 & 0.333284 & 1928258.797 & $+0.15770$ & -- & -- \\
\hline
2 & 160 & 58239 & 334809.1538 & 0.139157 & 57675.051 & $-0.00968$ & 93183 & 46590 \\
2 & 180 & 74192 & 433460.4919 & 0.136930 & 73475.274 & $-0.00966$ & 118708 & 59354 \\
2 & 200 & 92116 & 545868.2190 & 0.135001 & 91228.178 & $-0.00964$ & 147386 & 73692 \\
2 & 220 & 112018 & 672247.3736 & 0.133306 & 110941.295 & $-0.00961$ & -- & -- \\
2 & 250 & 145594 & 888433.1920 & 0.131102 & 144200.860 & $-0.00957$ & -- & -- \\
2 & 280 & 183655 & 1137075.8483 & 0.129212 & 181905.785 & $-0.00952$ & -- & -- \\
2 & 300 & 211529 & 1321128.5853 & 0.128090 & 209520.706 & $-0.00949$ & -- & -- \\
2 & 350 & 289997 & 1846288.9361 & 0.125656 & 287265.236 & $-0.00942$ & -- & -- \\
2 & 401 & 383013 & 2479252.0945 & 0.123590 & 379430.733 & $-0.00935$ & -- & -- \\
2 & 450 & 484778 & 3181621.6066 & 0.121895 & 480274.045 & $-0.00929$ & -- & -- \\
2 & 500 & 601189 & 3995002.3721 & 0.120388 & 595639.350 & $-0.00923$ & -- & -- \\
2 & 600 & 872259 & 5919943.2731 & 0.117874 & 864299.343 & $-0.00913$ & -- & -- \\
\hline
3 & 60 & 21175 & 919226.7397 & 0.018429 & 21178.048 & $+0.00014$ & 33882 & 16940 \\
3 & 70 & 29447 & 1462871.7509 & 0.016104 & 29449.642 & $+0.00009$ & 47116 & -- \\
3 & 80 & 39161 & 2187198.6384 & 0.014324 & 39164.223 & $+0.00008$ & 62659 & -- \\
3 & 90 & 50339 & 3118130.1632 & 0.012915 & 50342.246 & $+0.00006$ & -- & -- \\
3 & 100 & 62998 & 4281589.0863 & 0.011771 & 63001.731 & $+0.00006$ & -- & -- \\
3 & 110 & 77155 & 5703498.1687 & 0.010822 & 77158.783 & $+0.00005$ & -- & -- \\
3 & 120 & 92824 & 7409780.1714 & 0.010022 & 92827.959 & $+0.00004$ & -- & -- \\
3 & 130 & 110018 & 9426357.8555 & 0.009337 & 110022.540 & $+0.00004$ & -- & -- \\
3 & 140 & 128750 & 11779153.9820 & 0.008744 & 128754.741 & $+0.00004$ & -- & -- \\
3 & 150 & 149031 & 14494091.3118 & 0.008226 & 149035.865 & $+0.00003$ & -- & -- \\
3 & 160 & 170871 & 17597092.6061 & 0.007768 & 170876.432 & $+0.00003$ & -- & -- \\
\end{tabular}

\caption{Selected certified modes of the extended ranges (conventions of Tables~\ref{tab:m1}--\ref{tab:m3}; the rows $N=600$, $160$, $60$ repeat the last rows of those tables).}\label{tab:mX}
\end{table}

\begin{remark}[margins, and agreement with the floating-point values]
In the C certificates with a unique maximiser the gap between the lower bound for $2^P\sigma_{t^*}$ and the upper bounds at $t^*\pm1$ exceeds the width of the enclosure by a factor of at least \XMinMargin; the width $F^+_{t^*}-F^-_{t^*}$ never exceeds $\XMaxWidth$ units of $2^{-112}$. The longest run has $T=\XMaxT$ time steps and the largest reduced lattice $\XMaxSites$ sites; the floating phase never lasts more than $\XMaxTsw$ steps. The certified modes coincide with all \XFloatStudyAgree\ corner-to-corner values of the computations of Sections~5--6 of the article for $q\in\{4/5,1/2,1\}$ that fall in the sets $\mathcal N_{\rm C}$ (\XFloatStudyDisagree\ disagreements), including its largest cases $d=1$, $N=1000$; $d=2$, $N=401$; $d=3$, $N=100$.
% src: data/msc_rigorous_certified/SUMMARY_EXT.json (keys modes, comparison_with_float_study_modes); computations of Sections~5--6 of the article: data/msc_modes/discrete_modes.jsonl
\end{remark}

\begin{corollary}[the mode scales like $1/q$ up to one time step; extended]\label{cor:qscalingX}
With $\Delta_N=\tfrac45\,t^*(\tfrac45)-\tfrac12\,t^*(\tfrac12)$ as in Corollary~\ref{cor:qscaling}:
$\XQsLoOne\le\Delta_N\le\XQsHiOne$ for $d=1$, $2\le N\le\XQsNOne$; $\ \XQsLoTwo\le\Delta_N\le\XQsHiTwo$ for $d=2$, $2\le N\le\XQsNTwo$; $\ \XQsLoThree\le\Delta_N\le\XQsHiThree$ for $d=3$, $2\le N\le\XQsNThree$; all six bounds are attained.
% src: data/msc_rigorous_certified/SUMMARY_EXT.json (key q_scaling)
\end{corollary}
\begin{proof}
Exact arithmetic on the integers of Theorem~\ref{thm:modesX} (\texttt{summarise\_ext.py}, re-checked by \texttt{audit\_ext.py}).
\end{proof}

\begin{theorem}[accuracy of the closed form on the extended ranges; computer-assisted, trusted bases (T3), (T5)]\label{thm:cfX}
Let $t^*$ be the certified mode of Theorem~\ref{thm:modesX} and $t_{\rm cf}$ as in Definition~\ref{def:cf}.
\begin{itemize}
\item[(a)] $d=2$: for $q=4/5$ and all $10\le N\le\XEpsNTwoF$\XAlsoTwoF: $\ |t^*-t_{\rm cf}|\le\XEpsTwoF\;t^*$ (worst case $N=\XEpsAtTwoF$); for $q=1/2$ and all $10\le N\le\XEpsNTwoH$\XAlsoTwoH: $\ |t^*-t_{\rm cf}|\le\XEpsTwoH\;t^*$ (worst case $N=\XEpsAtTwoH$).
\item[(b)] $d=3$: for $q=4/5$ and all $10\le N\le\XEpsNThreeF$\XAlsoThreeF: $\ |t^*-t_{\rm cf}|\le\XEpsThreeF\;t^*$; for $q=1/2$ and all $10\le N\le\XEpsNThreeH$\XAlsoThreeH: $\ |t^*-t_{\rm cf}|\le\XEpsThreeH\;t^*$ (worst case $N=\XEpsAtThreeF$, resp.\ $N=\XEpsAtThreeH$).
\item[(c)] $d=1$: for $q=4/5$ and all $10\le N\le\XEpsNOneF$\XAlsoOneF: $\ \XEpsLoOneF\;t^*\le t_{\rm cf}-t^*\le\XEpsOneF\;t^*$; for $q=1/2$ and all $10\le N\le\XEpsNOneH$\XAlsoOneH: $\ \XEpsLoOneH\;t^*\le t_{\rm cf}-t^*\le\XEpsOneH\;t^*$.
\end{itemize}
Bounds on sub-ranges are given in Table~\ref{tab:cfX}; the enclosure of the relative error for every single $N$ is in \texttt{certs/closedform\_c128\_d$d$\_q$q$.jsonl}.
% src: data/msc_rigorous_certified/closedform_summary_c128.json; data/msc_rigorous_certified/closedform_c128_d*_q*.jsonl; re-checked with mpmath.iv in code/msc_rigorous_certified_c128/audit_ext.py
\end{theorem}
\begin{proof}
As for Theorem~\ref{thm:cf} (\texttt{scripts/closed\_form.py c128}: Arb at 256 bits from the rational enclosures of $\tau_N$ of Lemma~\ref{lem:mfpt}, which were computed for every $N$ of the sets $\mathcal N_{\rm C}(d,4/5)$; all have relative width below $10^{-28}$). Every bound of the statement and every row of Table~\ref{tab:cfX} was re-evaluated with \texttt{mpmath.iv} for every $N$ of its range by \texttt{audit\_ext.py}.
\end{proof}

\begin{table}[p]\centering\footnotesize
% generated by scripts/c128/summarise_ext.py -- do not edit
% src: data/msc_rigorous_certified/closedform_summary_c128.json (from data/msc_rigorous_certified/closedform_c128_d*_q*.jsonl, written by code/msc_rigorous_certified/closed_form.py c128)
\begin{tabular}{cc|c|c|c|c}
$d$ & $q$ & range of $N$ & bound on $|t_{\rm cf}-t^*|/t^*$ & worst $N$ & interval containing $(t_{\rm cf}-t^*)/t^*$ \\ \hline
1 & $1/2$ & $10\le N\le 1000$ & $0.228187$ & 10 & $[0.157934,\ 0.228187]$ \\
1 & $1/2$ & $20\le N\le 1000$ & $0.184401$ & 21 & $[0.157934,\ 0.184401]$ \\
1 & $1/2$ & $35\le N\le 1000$ & $0.171773$ & 36 & $[0.157934,\ 0.171773]$ \\
1 & $1/2$ & $50\le N\le 1000$ & $0.167141$ & 51 & $[0.157934,\ 0.167141]$ \\
1 & $1/2$ & $60\le N\le 1000$ & $0.165552$ & 60 & $[0.157934,\ 0.165552]$ \\
1 & $1/2$ & $100\le N\le 1000$ & $0.162171$ & 102 & $[0.157934,\ 0.162171]$ \\
1 & $1/2$ & $160\le N\le 1000$ & $0.160388$ & 161 & $[0.157934,\ 0.160388]$ \\
1 & $1/2$ & $200\le N\le 1000$ & $0.159792$ & 200 & $[0.157934,\ 0.159792]$ \\
1 & $1/2$ & $300\le N\le 1000$ & $0.159009$ & 300 & $[0.157934,\ 0.159009]$ \\
1 & $1/2$ & $600\le N\le 1000$ & $0.158243$ & 600 & $[0.157934,\ 0.158243]$ \\
1 & $4/5$ & $10\le N\le 1500$ & $0.225045$ & 11 & $[0.157780,\ 0.225045]$ \\
1 & $4/5$ & $20\le N\le 1500$ & $0.186954$ & 21 & $[0.157780,\ 0.186954]$ \\
1 & $4/5$ & $35\le N\le 1500$ & $0.172733$ & 35 & $[0.157780,\ 0.172733]$ \\
1 & $4/5$ & $50\le N\le 1500$ & $0.167595$ & 50 & $[0.157780,\ 0.167595]$ \\
1 & $4/5$ & $60\le N\le 1500$ & $0.165848$ & 60 & $[0.157780,\ 0.165848]$ \\
1 & $4/5$ & $100\le N\le 1500$ & $0.162246$ & 101 & $[0.157780,\ 0.162246]$ \\
1 & $4/5$ & $160\le N\le 1500$ & $0.160415$ & 161 & $[0.157780,\ 0.160415]$ \\
1 & $4/5$ & $200\le N\le 1500$ & $0.159836$ & 200 & $[0.157780,\ 0.159836]$ \\
1 & $4/5$ & $300\le N\le 1500$ & $0.159020$ & 300 & $[0.157780,\ 0.159020]$ \\
1 & $4/5$ & $600\le N\le 1500$ & $0.158247$ & 600 & $[0.157780,\ 0.158247]$ \\
1 & $4/5$ & $N=2000$ & $0.157705$ & 2000 & $[0.157704,\ 0.157705]$ \\
2 & $1/2$ & $10\le N\le 200$ & $0.009704$ & 133 & $[-0.009704,\ 0.006710]$ \\
2 & $1/2$ & $20\le N\le 200$ & $0.009704$ & 133 & $[-0.009704,\ -0.004819]$ \\
2 & $1/2$ & $35\le N\le 200$ & $0.009704$ & 133 & $[-0.009704,\ -0.007947]$ \\
2 & $1/2$ & $50\le N\le 200$ & $0.009704$ & 133 & $[-0.009704,\ -0.008933]$ \\
2 & $1/2$ & $60\le N\le 200$ & $0.009704$ & 133 & $[-0.009704,\ -0.009238]$ \\
2 & $1/2$ & $100\le N\le 200$ & $0.009704$ & 133 & $[-0.009704,\ -0.009634]$ \\
2 & $1/2$ & $160\le N\le 200$ & $0.009690$ & 160 & $[-0.009690,\ -0.009640]$ \\
2 & $4/5$ & $10\le N\le 301$ & $0.009698$ & 137 & $[-0.009698,\ 0.005978]$ \\
2 & $4/5$ & $20\le N\le 301$ & $0.009698$ & 137 & $[-0.009698,\ -0.004006]$ \\
2 & $4/5$ & $35\le N\le 301$ & $0.009698$ & 137 & $[-0.009698,\ -0.007897]$ \\
2 & $4/5$ & $50\le N\le 301$ & $0.009698$ & 137 & $[-0.009698,\ -0.008910]$ \\
2 & $4/5$ & $60\le N\le 301$ & $0.009698$ & 137 & $[-0.009698,\ -0.009204]$ \\
2 & $4/5$ & $100\le N\le 301$ & $0.009698$ & 137 & $[-0.009698,\ -0.009491]$ \\
2 & $4/5$ & $160\le N\le 301$ & $0.009687$ & 163 & $[-0.009687,\ -0.009491]$ \\
2 & $4/5$ & $200\le N\le 301$ & $0.009639$ & 200 & $[-0.009639,\ -0.009491]$ \\
2 & $4/5$ & $300\le N\le 301$ & $0.009495$ & 300 & $[-0.009495,\ -0.009491]$ \\
2 & $4/5$ & $N=350$ & $0.009420$ & 350 & $[-0.009420,\ -0.009419]$ \\
2 & $4/5$ & $N=401$ & $0.009353$ & 401 & $[-0.009353,\ -0.009352]$ \\
2 & $4/5$ & $N=450$ & $0.009291$ & 450 & $[-0.009291,\ -0.009290]$ \\
2 & $4/5$ & $N=500$ & $0.009232$ & 500 & $[-0.009232,\ -0.009231]$ \\
2 & $4/5$ & $N=600$ & $0.009126$ & 600 & $[-0.009126,\ -0.009125]$ \\
3 & $1/2$ & $10\le N\le 80$ & $0.001562$ & 10 & $[0.000052,\ 0.001562]$ \\
3 & $1/2$ & $20\le N\le 80$ & $0.000386$ & 20 & $[0.000052,\ 0.000386]$ \\
3 & $1/2$ & $35\le N\le 80$ & $0.000182$ & 35 & $[0.000052,\ 0.000182]$ \\
3 & $1/2$ & $50\le N\le 80$ & $0.000111$ & 51 & $[0.000052,\ 0.000111]$ \\
3 & $1/2$ & $60\le N\le 80$ & $0.000085$ & 60 & $[0.000052,\ 0.000085]$ \\
3 & $4/5$ & $10\le N\le 120$ & $0.002749$ & 10 & $[0.000039,\ 0.002749]$ \\
3 & $4/5$ & $20\le N\le 120$ & $0.000945$ & 21 & $[0.000039,\ 0.000945]$ \\
3 & $4/5$ & $35\le N\le 120$ & $0.000249$ & 37 & $[0.000039,\ 0.000249]$ \\
3 & $4/5$ & $50\le N\le 120$ & $0.000195$ & 50 & $[0.000039,\ 0.000195]$ \\
3 & $4/5$ & $60\le N\le 120$ & $0.000144$ & 60 & $[0.000039,\ 0.000144]$ \\
3 & $4/5$ & $100\le N\le 120$ & $0.000060$ & 100 & $[0.000039,\ 0.000060]$ \\
3 & $4/5$ & $N=130$ & $0.000042$ & 130 & $[0.000041,\ 0.000042]$ \\
3 & $4/5$ & $N=140$ & $0.000037$ & 140 & $[0.000036,\ 0.000037]$ \\
3 & $4/5$ & $N=150$ & $0.000033$ & 150 & $[0.000032,\ 0.000033]$ \\
3 & $4/5$ & $N=160$ & $0.000032$ & 160 & $[0.000031,\ 0.000032]$ \\
\end{tabular}

\caption{Certified bounds for the relative error of the closed form on sub-ranges of the extended ranges (conventions of Table~\ref{tab:cf}).}\label{tab:cfX}
\end{table}

\begin{remark}\label{rem:cfX}
The extension confirms the picture of Remark~\ref{rem:cfshape} with certified numbers (all taken from Table~\ref{tab:cfX}). In two dimensions the relative error is negative for every $N\ge14$ of the ranges (both values of $q$); its modulus has a flat maximum, $\XEpsTwoF$ at $N=\XEpsAtTwoF$ for $q=4/5$, and decreases slowly afterwards, apart from fluctuations of order $1/t^*$ caused by the integer rounding of the mode (rows $160\le N\le\XEpsNTwoF$ and $200\le N\le\XEpsNTwoF$ and the isolated values of $N$ in the table). In three dimensions the error is positive for every $N\ge3$ of the ranges at $q=4/5$ and for every $N\ne3$ at $q=1/2$; it falls to $3.17888\times10^{-5}$ (rounded) at $N=160$ ($q=4/5$), though not monotonically, for the same reason. In one dimension $t_{\rm cf}/t^*-1$ lies in $[0.157704,0.157705]$ at $N=2000$; the value suggested by the continuum limit ($X_N\to\pi^2/2$ and $t^*/\mathbb E\,T\to\kappa_1$, neither of which is proved in this note) is $\ln(\pi^2)/\bigl((\pi^2/2+1)\kappa_1\bigr)-1=0.157475\ldots$ The closed form is not an approximation of the one-dimensional mode.
% src: data/msc_rigorous_certified/remark_checks.json (keys closed_form_shape, remark_cfX_continuum_value); data/msc_rigorous_certified/closedform_c128_d1_q4-5.jsonl (N = 2000), data/msc_rigorous_certified/closedform_c128_d3_q4-5.jsonl (N = 160)
\end{remark}

\begin{theorem}[the one-dimensional mode on the extended range; computer-assisted, trusted bases (T3), (T5)]\label{thm:1dlawX}
Let $d=1$, $\kappa_1$ as in Theorem~\ref{thm:tau} and $y_N(q)=\kappa_1(N-\frac12)^2/q$.
\begin{itemize}
\item[(a)] For $q=4/5$ and every $10\le N\le\XNmaxOneF$\XAlsoOneF: $\ \XDeltaLoF<t^*-y_N<\XDeltaHiF$, hence $t^*=\lfloor y_N\XDeltaHiF\rfloor$.
\item[(b)] For $q=1/2$ and every $10\le N\le\XNmaxOneH$\XAlsoOneH: $\ \XDeltaLoH<t^*-y_N<\XDeltaHiH$, hence $t^*=\lfloor y_N\XDeltaHiH\rfloor$.
\item[(c)] On the contiguous ranges, $\bigl|t^*/\mathbb E\,T-\kappa_1\bigr|\le c/(N(N-1))$ with $c=\XRatioConstF$ for $q=4/5$ and $c=\XRatioConstH$ for $q=1/2$.
% src: data/msc_rigorous_certified/SUMMARY_EXT.json (key law_1d), from data/msc_rigorous_certified/closedform_c128_d1_q*.jsonl and data/msc_rigorous_certified/constants.json
\end{itemize}
\end{theorem}
\begin{proof}
Exact rational arithmetic on the integers of Theorem~\ref{thm:modesX} and the rational enclosure of $\kappa_1$ (\texttt{closed\_form.py c128}, \texttt{summarise\_ext.py}; re-checked for every $N$, with both endpoints of the enclosure of $\kappa_1$, by \texttt{audit\_ext.py}).
\end{proof}
